\section{总结}

这个项目我是从2015年4月20日开始，基本代码完成于5月19日。最后到小学期，完成的这个文档。
整个过程是很有意思的，也成为了我在GitHub上提交最频繁的一段记录。
当然，还是一直在参考书\cite{osask}中的内容再做。
书的内容是中文的，但是给出的随书源代码是日文注释，不过正好会日语，所以对源代码的理解和重写也就轻松了许多。
许多函数名和变量名，实际上也都是日语的一些东西，看懂了是挺有意思的，看得出作者的用心。
比如Haribote这个词就是{\ja 「張りぼて」}，意思是“纸糊的戏剧用小道具”。
也就是说，这个操作系统，就像纸糊的一样，只是个玩具，看起来还不错，其实是空的。
所以，我们也可以看到，这个操作系统，真的就像玩具一样，可以玩玩，但并不能实用。
比如它不能真正地读写文件，没有虚拟内存，不能连接网络等。

在现代操作系统中，进程、存储、I/O、文件、网络，显然是必要的五个元素。
这里呢，实际上只是对进程的实现，是相对比较丰富的。
从代码中我们可以看到，这个操作系统，是分时多任务的。
也就是说，我们为每个进程分配一个时间片，然后切换上下文来运行在队列中的进程。
这里我们也加了一个优先级，来调度不同等级的任务。
查找算法，也都是非常简单的线性算法。
因为任务很少，也没有明显的性能问题。

对于存储器的管理，实际上更简单了。
我们维护一个空闲内存块的链表，然后分配和释放都在这个链表上进行。
算法也是线性的。

I/O存取，其实这里就是对系统中断的处理了。
主要是鼠标和键盘中断，以及屏幕的绘制。

屏幕绘制，专门写了一个窗口和图层管理器。
这个算是代码量最大的一部分了，所以看起来这个操作系统代码很多，大部分都不是核心。
写屏幕很有意思，因为能够看到花花绿绿的结果，很开心，有成就感。
但实际上这些东西并不是操作系统的核心。

当然完成了这些东西，提供了一些API供应用程序使用，精力也就比较有限了。
所以文件系统没有实现，只是简单地在汇编阶段把软盘（而且只支持FAT12格式的软盘）内容全部读到内存中，再去操作内存。
同样，网络也没有实现，下学期学《计算机网络》的时候，倒不妨写一个TCP/IP玩玩。

在此期间我又读了读MINIX\cite{minix3}的源代码，这是一个分模块的操作系统实现，也从中学到了很多东西。
不过还没有去读Linux的内核源代码，这是计划了。

其实这个学期对我影响最深的是CS:APP——《深入理解计算机系统》\cite{csappv2}。
花了两个月的时间读完了。
这本书从头到尾把计算机硬件到软件讲了一遍，是很好的计算机入门书。
而且也正式从这本书中学到了汇编，对C的理解更深了一些。
不过学习的教学环境还没有这种课程，我对于这本书的理解也只是浅尝辄止，后面的作业其实没有做多少。
参考文献列在最后，\LaTeX 用来排版以及{\sc Bib}\TeX 用来列参考文献，倒是挺方便的。

操作系统原理性的东西，其实在上个世纪80年代已经全部搞明白了。
作为一个科班出身的程序员，能写OS也不是什么稀奇事了。
但是，玩具好做，真正能够把操作系统上的API设计得好，能够构建出一系列应用程序，
是一项非常浩大的工程。从内核构建到虚拟机，再到用户界面和应用程序，
只能说只有微软完成了这个完整地流程。
就算是Mac OS X和iOS，也是基于UNIX实现改过来的，说白了不过加个壳。
而Linux至今在桌面领域一团糟，衍生而来的Android虽然在移动领域二分天下，
不过也不是Google独立开发出来的，还面临着Java的版权问题。
所以这么一个庞大的生态系统，可以说难于上青天。当然现在上天倒是很随意的事情了。

话说回来，这个操作系统，还只能是玩具。
作为课程设计，这个可能有点太大了，但我觉得这一轮下来也是值得的。
这回课程设计，能够借机完成一个操作系统，也是我一直以来的愿望。
不过其实对这个玩具并不满意，除了可能GUI设计的扁平化了一些，看起来更像现代操作系统以外，
真是没有一点可爱的地方。本来还想实现汉字显示，结果折腾了半个月没弄好，
还好一直用git做版本控制，就回滚到能用的版本提交了这个版本。
如果有时间，我倒是想写一个UNIX兼容（POSIX）的内核，而不是现在这个自有API的东西。
当然这是后话了。

更多的探索还在后面，这个课程设计的完成也只是更大目标的开始。
要学的东西还有很多，所以就去更多地学习和实践了！

\vfill

\centering
\includegraphics[width=12cm]{image/laala.png}

\vspace*{1cm}
